\documentclass[10pt,a4paper]{amsart}
\usepackage[margin=19mm]{geometry}
\usepackage{amsmath,amssymb,mathtools}
\usepackage[scr=rsfs,cal=euler]{mathalfa}
\usepackage{xfrac}
\usepackage[unicode,hidelinks,bookmarks=false]{hyperref}
\newcommand{\ie}{i.e.\ }
\newcommand{\cf}{cf.\ }
\newcommand{\resp}{respectively\ }
\newcommand{\Subsection}{\subsection}
\providecommand{\nobreakdash}{\nobreak\hbox{-}}
\setlength{\emergencystretch}{2em}
\multlinegap0pt
\usepackage{amssymb}
\usepackage{xcolor}
\usepackage[shortlabels]{enumitem}
\newtheorem{lemma}{Lemma}
\newcommand{\azul}[1]{\textcolor{black}{#1}}
\title{Factorization patterns and fields of definition of $\ell$-torsion points on the Jacobians of genus 3 hyperelliptic curves}
\author{Amalia Pizarro-Madariaga, Edgardo Riquelme}
\date{Source: arXiv:2609.03171v1 (CC BY 4.0)}
\begin{document}
\maketitle

\noindent\emph{Source and licence.} Factorization patterns and fields of definition of $\ell$-torsion points on the Jacobians of genus 3 hyperelliptic curves. Version arXiv:2609.03171v1 (CC BY 4.0), distributed by arXiv under the Creative Commons Attribution 4.0 International licence, http://creativecommons.org/licenses/by/4.0/, as recorded in the arXiv licence metadata for this exact version and shown on the abstract page https://arxiv.org/abs/2609.03171v1. Full licence notice: \texttt{licenses/arxiv-2609.03171-CC-BY-4.0.txt}.\par\smallskip

\noindent\textit{Editorial scope.} Counted unit: the article's lemma lemma:orbitlayers (source0.tex lines 349--357). The article's complete original proof of it is delivered with it (source0.tex lines 358--370, one \texttt{proof} environment of 13 lines). No mathematics is written, repaired or re-derived: the statement and the proof are the article's own lines, in the article's own notation, and no retained line is substituted or re-flowed.  No further source-local context is needed: every cross-reference of every retained line resolves inside the two delivered spans.  Nothing was omitted from any retained span, so the package carries no omission record.  No retained line contains a citation, so the wrapper carries no bibliography and no bibliography entry.  No retained line contains a figure environment, an image-inclusion command or drawing code, so no image is delivered and the package declares no asset dependency.  Editorial formatting only, in addition: an AMS \texttt{amsart} preamble that restates the article's own declarations and macros used by the retained lines, namely the environments \texttt{lemma} on separate counters, together with the standard packages those lines need.  The retained environments print in this document as Lemma 1 in order of appearance, whereas the article prints the same items with the numbering of its own full document.  The complete native source of the article as distributed in the arXiv e-print for this exact version is bundled unchanged as \texttt{sources/209271/source0.tex}.
\begin{lemma}\label{lemma:orbitlayers}
\azul{Let $A_d$ be the block associated with an irreducible polynomial
$p(x)$ of degree $d$, and let $n$ be its length.
For each $1\le k\le n$, the number of elements annihilated by
$p(x)^k$ but not by $p(x)^{k-1}$ is
\[
\ell^{d(k-1)}(\ell^d-1).
\]}
\end{lemma}

\begin{proof}
The subspace annihilated by $p(x)^k$ has dimension $kd$ and hence
contains $\ell^{kd}-1$ nonzero elements.
Similarly, the subspace annihilated by $p(x)^{k-1}$ has dimension
$(k-1)d$ and contains $\ell^{(k-1)d}-1$ nonzero elements.
Therefore the number of elements annihilated by $p(x)^k$ but not by
$p(x)^{k-1}$ is
\[
(\ell^{kd}-1)-(\ell^{(k-1)d}-1)
=
\ell^{(k-1)d}(\ell^d-1).
\]
\end{proof}

\end{document}
